% Author: Grayson Orr
% Course: ID511001: Programming 2

\documentclass{article}
\author{}

\usepackage{graphicx}
\usepackage{wrapfig}
\usepackage{enumerate}
\usepackage{hyperref}
\usepackage[margin = 2.25cm]{geometry}
\usepackage[table]{xcolor}
\usepackage{soul}
\usepackage{fancyhdr}
\usepackage{fancyvrb}
\hypersetup{
  colorlinks = true,
  urlcolor = blue
}
\setlength\parindent{0pt}
\pagestyle{fancy}
\fancyhf{}
\rhead{College of Engineering, Construction and Living Sciences\\Bachelor of Information Technology}
\lfoot{Project 2: Pong\\Version 3, Semester Two, 2023}
\rfoot{\thepage}
 
\begin{document}

\begin{figure}
    \centering
    \includegraphics[width=50mm]{../../resources/img/logo.png}
\end{figure}

\title{College of Engineering, Construction and Living Sciences\\Bachelor of Information Technology\\ID511001: Programming 2\\Level 5, Credits 15\\\textbf{Project 2: Pong}}
\date{}
\maketitle

\section*{Assessment Overview}
In this assessment, you will design and develop a pong \textbf{Windows Forms App} using \textbf{C\#}.

\section*{Learning Outcomes}
At the successful completion of this course, learners will be able to:
\begin{enumerate}
    \item Build interactive, event-driven GUI applications using pre-built components.
    \item Declare and implement user-defined classes using encapsulation, inheritance and polymorphism.
\end{enumerate} 

\section*{Assessments}
\renewcommand{\arraystretch}{1.5}
\begin{tabular}{|c|c|c|c|}
	\hline
	\textbf{Assessment}                                 & \textbf{Weighting} & \textbf{Due Date}            & \textbf{Learning Outcomes} \\ \hline
	\small Project 1: Student Management System  & \small 35\%        & \small 22-09-2023 (Friday at 4.59 PM)   & \small 1 and 2                   \\ \hline
	\small Project 2: Pong & \small 25\%        & \small 10-11-2023 (Friday at 04.59 PM)  & \small 1 and 2                   \\ \hline
	\small Theory Examination                        & \small 30\%        & \small 15-11-2023 (Wednesday at 12.10 PM)  & \small 1 and 2                   \\ \hline
	\small Classroom Task: Unit Testing                       & \small 10\%        & \small 22-09-2023 (Friday at 4.59 PM)  & \small 1 and 2                   \\ \hline
\end{tabular} 

\section*{Conditions of Assessment}
You will complete this assessment during your learner-managed time. However, there will be time during class to discuss the requirements and your progress on this assessment. This assessment will need to be completed by \textbf{Wednesday, 14 June 2023} at \textbf{4.59 PM}.

\section*{Pass Criteria}
This assessment is criterion-referenced (CRA) with a cumulative pass mark of \textbf{50\%} over all assessments in \textbf{ID511001: Programming 2}.

\section*{Authenticity}
All parts of your submitted assessment \textbf{must} be completely your work. If you use code snippets from \textbf{GitHub}, \textbf{StackOverflow} or other online resources, you \textbf{must} reference it appropriately using \textbf{APA 7th edition}. Provide your references in the \textbf{README.md} file in your repository. Failure to do this will result in a mark of \textbf{zero} for this assessment.

\section*{Policy on Submissions, Extensions, Resubmissions and Resits}
The school's process concerning submissions, extensions, resubmissions and resits complies with \textbf{Otago Polytechnic | Te Pūkenga} policies. Learners can view policies on the \textbf{Otago Polytechnic | Te Pūkenga} website located at \href{https://www.op.ac.nz/about-us/governance-and-management/policies}{https://www.op.ac.nz/about-us/governance-and-management/policies}.

\section*{Submission}
You \textbf{must} submit all application files via \textbf{GitHub Classroom}. Here is the URL to the repository you will use for your submission – \href{https://classroom.github.com/a/eFe1Oh97}{https://classroom.github.com/a/eFe1Oh97}.  Create a \textbf{.gitignore} and add the ignored files in this resource - \href{https://raw.githubusercontent.com/github/gitignore/main/VisualStudio.gitignore}{https://raw.githubusercontent.com/github/gitignore/main/VisualStudio.gitignore}. The latest application files in the \textbf{master} or \textbf{main} branch will be used to mark against the \textbf{Functionality} criterion. Please test before you submit. Partial marks \textbf{will not} be given for incomplete functionality. Late submissions will incur a \textbf{10\% penalty per day}, rolling over at \textbf{5:00 PM}.

\section*{Extensions}
Familiarise yourself with the assessment due date. Contact the course lecturer before the due date if you need an extension. If you require more than a week's extension, you will need to provide a medical certificate or support letter from your manager.

\section*{Resubmissions}
Learners may be requested to resubmit an assessment following a rework of part/s of the original assessment. Resubmissions are to be completed within a negotiable short time frame and usually \textbf{must} be completed within the timing of the course to which the assessment relates. Resubmissions will be available to learners who have made a genuine attempt at the first assessment opportunity and achieved a \textbf{D grade (40-49\%)}. The maximum grade awarded for resubmission will be \textbf{C-}.

\section*{Resits}
Resits and reassessments \textbf{are not} applicable in \textbf{ID511001: Programming 2}.

\section*{Instructions}
You will need to submit an application and documentation that meet the following requirements:\\

\subsection*{Functionality - Learning Outcomes 1 and 2 (50\%)}
    \item The application needs to open without code or file structure modification in \textbf{Visual Studio}.
    \item The game needs to be driven by one \textbf{Timer} and begins when the user presses the \textbf{space bar} key.
    \item The ball and two paddles need to be created using the \textbf{Graphics} class. 
    \item The ball needs to bounce/collide off the top and bottom of the screen, and paddles.
    \item The paddles needs to move vertically but not exceed the top and bottom of the screen.
    \item The user controls the left paddle via the \textbf{up arrow} and \textbf{down arrow} keys. The computer controls the right paddle. It is acceptable for the right paddle to follow the ball's position. However, other solutions are encouraged.
    \item Display the user and computer's score using the \textbf{DrawString} method.
    \item A scoring system. When the ball collides with the left and right-hand side of the screen, one point is given to either the user or computer. The game is over when either score is 10. 
    \item A highscore system. When the game is over, appropriate feedback needs to be displayed to the user, i.e., \textbf{"You win!"} or \textbf{"You lose!"}, the user and computer's scores are saved, i.e., written to a text file. Read the scores from the text file and display the last five to the user.
    \item Double buffering to prevent the ball, paddles and scores from flickering.
    \item Using the \textbf{SoundPlayer} class, play a sound when:
    \begin{itemize}
        \item The ball bounces off the paddle, and top and bottom of the screen.
        \item The user wins and loses.
    \end{itemize}
    \item Restart and pause a game via the \textbf{R} and \textbf{P} keys.
    \item Randomise the colour of the ball and paddles.
\end{itemize}

\subsection*{Code Elegance - Learning Outcomes 1 and 2 (40\%)}
\begin{itemize}
    \item A \textbf{Visual Studio} \textbf{.gitignore} file is used. 
    \item Appropriate naming of files, variables, methods and classes.
    \item Idiomatic use of control flow, data structures and in-built functions.
    \item Efficient algorithmic approach.
    \item Sufficient modularity.
    \item Each files needs to have a header comment located at the top of the file.
    \item In-line comments where required. It should be for code that needs further explanation.
    \item Formatted code.
    \item No dead or unused code.
\end{itemize}

\subsection*{Documentation and Git Usage - Learning Outcomes 1 and 2 (10\%)}
\begin{itemize}
    \item Provide the following in your repository \textbf{README.md} file:
    \begin{itemize}
        \item A class diagram of your application.
        \item If applicable, known bugs.
    \end{itemize}
    \item Commit messages reflect the context of each functional requirement change.
\end{itemize}

\subsection*{Additional Information}
\begin{itemize}
    \item An exemplar is available in the \textbf{assessment} directory of the \textbf{course materials} repository.
    \item You may add additional classes and methods.
    \item When the user presses a key, i.e., \textbf{up arrow} or \textbf{down arrow}, a \textbf{KeyDown} event is generated. For the \textbf{Form1's} \textbf{KeyDown} event, the method signature is:
    \begin{Verbatim}[tabsize=2]
		private void Form1_KeyDown(object sender, KeyEventArgs e) {}
	\end{Verbatim}
	The argument you will be interested in is \textbf{KeyEventArgs e} which is the value of the pressed key. The \textbf{arrow} key values are \textbf{Keys.Left}, \textbf{Keys.Right}, \textbf{Keys.Up} and \textbf{Keys.Down}. In the \textbf{Form1\_KeyDown} method, you can use a \textbf{switch} statement. For example:
	\begin{Verbatim}[tabsize=2]
		switch (e.KeyCode) 
		{
			case Keys.Left:
				// Do something
				break;
			case Keys.Right:
				// Do something
				break;
			case Keys.Up:
				// Do something
				break;
			case Keys.Down:
				// Do something
				break;
			default:
				// Do something
				break;
		}
	\end{Verbatim}
	\textbf{Note:} The \textbf{Form1's} \textbf{KeyPreview} event needs to be set to \textbf{True}. Otherwise, \textbf{Form1} will not respond to the \textbf{KeyDown} event.
    \item \textbf{Do not} rewrite your \textbf{Git} history. It is important that the course lecturer can see how you worked on your assessment over time
\end{itemize}

\end{document}